\chapter{使用说明与复习路线}
\section*{三遍阅读}
第一遍用第1—6章建立模型与分布图景，遇题先写样本空间或支撑集。第二遍完成第7—8章，把“估计量—分布—区间—拒绝域”连接起来。第三遍回到竞赛题和第9—10章，用不同方法重建结果。
\begin{center}\small
\begin{tabularx}{\textwidth}{@{}lXX@{}}\toprule
阶段&训练重点&自检标准\\\midrule
第1—2周&事件、条件概率、一元分布&能独立写事件并检查分布函数端点\\
第3—4周&多维分布、数字特征、极限&能画积分区域并说明独立性用途\\
第5—6周&抽样、估计、检验&能选枢轴量并解释自由度与尾部\\
第7—8周&回归、综合题、选学&能比较两法并检查假设改变的后果\\\bottomrule
\end{tabularx}\end{center}
\section*{盒子与题目层级}
定义和方法卡用青色；定理与例题用墨蓝；竞赛拓展和考研提示用金色；易错警示用赭色；测度论选学用浅墨色。星级从一星到三星表示本册内部相对难度。
每个知识节有3例与至少1张方法卡；章末8题依次包含基础、考研、竞赛、测度论选学各2题。点击题目后的答案页码可跳转，答案末可返回题目。
\section*{解题八问}
\begin{enumerate}
\item 样本点、随机变量或统计量分别是什么？
\item 参数空间和变量支撑是否依赖未知参数？
\item 独立、同分布、正态、有限矩分别用在哪里？
\item 是精确概率、上界，还是渐近近似？
\item 能否改求补事件、改用指示变量或先条件化？
\item 变换是否多对一，雅可比是否取绝对值？
\item 置信水平、显著性水平和分位数方向是否统一？
\item 结果能否通过总概率为1、量纲、极限或特例检查？
\end{enumerate}
\section*{统一约定}
所有分位数取左广义逆$q_u=\inf\{x:F(x)\ge u\}$；Gamma、指数均用率参数$\lambda$，若使用尺度参数$\eta$，两者关系为$\lambda=\frac{1}{\eta}$；iid表示独立同分布。测度论内容可在完成主线后阅读。
